\textbf{Proof.}
Choose a nonzero compactly supported function \(\varphi\) on \([-1,1]\) with \(\|\varphi\|_\infty=1\) and finite H\"older seminorm \([\varphi]_\beta\).  Since \(\mathcal I\) is a compact subinterval of \((0,1)\), for all sufficiently small \(h\) it contains \(m\asymp h^{-1}\) centers whose radius-\(h\) neighborhoods are pairwise disjoint.  Set
\[
 \tau_j(v)=\gamma h^\beta\varphi((v-v_j)/h),\quad j=1,\ldots,m,
 \qquad \tau_0(v)=0. \tag{1}
\]
Scaling gives \([\tau_j]_\beta\le\gamma[\varphi]_\beta\) and \(\|\tau_j\|_\infty=\gamma h^\beta\).  Choose \(\gamma>0\) small enough that \(\gamma[\varphi]_\beta\le L_\tau\), the range bound holds, and the divergence constant below is as small as required.  The established exact-code reduction then verifies every primitive law condition.  Since \(g_0=\bar g_n=0\) and \(0\in\mathcal G\), the code error is zero, so every pair lies in \(\mathcal Q_{\mathrm{spec},n}\) for the fixed positive admissible radius.

Disjoint support and \(\|\varphi\|_\infty=1\) imply, for distinct \(j,\ell\in\{0,\ldots,m\}\),
\[
 \|\tau_j-\tau_\ell\|_\infty\ge\gamma h^\beta. \tag{2}
\]
With \(h\asymp(\log n/n)^{1/(2\beta+1)}\), the right side is at least \(c r_n\).

Under the exact-code reduction, \(T\) has the same ancillary distribution under every hypothesis and, conditional on \(V\),
\[
 Y/T=\tau_j(V)+\epsilon,\qquad \epsilon\sim N(0,\sigma_k^2). \tag{3}
\]
The iid product-law divergence from the zero curve is therefore exactly
\[
 \operatorname{KL}(P_j^n,P_0^n)
 =\frac{n}{2\sigma_k^2}\int_0^1\tau_j(v)^2\,dv
 =\frac{\gamma^2\|\varphi\|_2^2}{2\sigma_k^2}nh^{2\beta+1}. \tag{4}
\]
Now \(nh^{2\beta+1}\asymp\log n\), while \(m\asymp h^{-1}\) implies \(\log m\sim(2\beta+1)^{-1}\log n\).  Decreasing the fixed \(\gamma\), if necessary, makes (4) at most \(c_0\log m\) for any prescribed sufficiently small \(c_0>0\).  This proves membership, separation, and the divergence claim. \(\square\)